\documentclass[conference,letterpaper]{IEEEtran}
\usepackage[T1]{fontenc}
\usepackage{amsmath,amssymb}
\usepackage{cite}
\usepackage{balance}
\usepackage[hidelinks]{hyperref}
\hypersetup{
  pdftitle={Adversarial Cooperation},
  pdfauthor={Santiago Restrepo},
  pdfsubject={Research proposal, September 2026},
  pdfkeywords={adversarial cooperation, applied cryptography, zero knowledge,
    multiparty computation, hardware},
  pdfdisplaydoctitle=true
}
\urlstyle{same}

\title{Adversarial Cooperation}
\author{\IEEEauthorblockN{Santiago Restrepo}
\IEEEauthorblockA{\href{mailto:contact@waajacu.com}{contact@waajacu.com}\quad
\url{https://waajacu.com/}\\[3pt]
\textit{Research proposal --- September 2026}}}
\date{}

\begin{document}
\maketitle

\begin{abstract}
Adversarial Cooperation investigates cryptographic protocols for classical
and modern problems in game theory without a trusted referee, and the
hardware that parties would use to execute them. The aim is to enable
cooperation between adversarial parties while preserving private information
and the strategic advantages that depend on it. Each construction will
define what must be learned to reach an agreed result and seek to reveal
nothing further under explicit assumptions. We seek resources and research
support to develop, analyze, implement, and demonstrate these protocols and
their hardware.
\end{abstract}

\begin{IEEEkeywords}
Applied cryptography, cooperation, zero knowledge, multiparty computation,
hardware.
\end{IEEEkeywords}

\section{The Cooperation Problem}
The central question is: \emph{how can we construct and demonstrate
cryptographic protocols that let adversarial parties cooperate without a
referee, while retaining the private information on which their strategic
advantage depends?} A party should be able to take part in a common game or
verify an agreed result without handing its cards, strategy, or other secrets
to an opponent or a trusted intermediary.

\emph{Author's intuition:} ``Playing poker without revealing the cards, or
playing chess without moving the pieces.'' In poker, a question is whether
players can verify a winning hand without disclosing its cards.

For chess, the intuition is that visible play may work against the parties'
interests by exposing information about their strategies. If both parties
fix deterministic move-selection algorithms, the proposed protocol would
privately compute the outcome of their matchup from an agreed initial
position. On completion, it would reveal only win, loss, or draw, seeking to
hide the algorithms and move sequence beyond what the result implies. The
agreed rules must specify draws and how illegal moves or exceeded computation
limits affect the result. The computation may simulate play internally; this
proposal does not assume a shortcut to determining the outcome.

The privacy objective is to reveal no additional private information beyond
the agreed outputs and public information, under the chosen adversary model.
The parties still exchange protocol messages, and the result itself may
reveal information. Preserving a strategic advantage therefore requires
identifying the secret it depends on and examining what the output, repeated
queries, and subsequent play disclose. Execution costs and aborts must also
be considered.

The program builds on commitments, zero-knowledge proofs, and
multiparty computation: established ways to fix private choices, prove predicates,
and evaluate joint functions under specified assumptions
\cite{naor1991,gmr1989,gmw1987}. Cryptographic replacement of a mediator in
strategic games is also prior work \cite{dodis2000}. This proposal claims no
new primitive or general solution to cooperation. It proposes carefully
specified case studies that connect these foundations to the decisions
parties face, and follow the resulting requirements into software and
hardware. Whether any particular composition constitutes a research
contribution must be established case by case.

\section{Scope, Assumptions, and Requirements}
Here, \emph{without a trusted authority} means that the selected task requires
no central party entrusted with all inputs or unilateral authority over its
result. It does not mean absence of assumptions. Each case study must state
its authentication mechanism, channel model, randomness requirements, setup
resources, corruption threshold, and reliance on local devices. A setup that
introduces a trusted dealer must be identified as changing the objective.

The proposed starting model has two parties and authenticated channels, with
at most one party statically corrupted by a probabilistic polynomial-time
adversary. The adversary can stop or delay communication. Each construction
must specify any additional channel or setup requirements and how they are
met. Computational privacy and correctness against malicious deviation are
initial research targets under these assumptions; they remain to be
established for the proposed implementations.

For each task, an ideal functionality will specify inputs, outputs, allowed
leakage, and the intended behavior of output release. Each party would be
required to assess the algorithm's privacy and correctness under the stated
adversary model and conclude that the evidence justifies its use. It would
then fabricate its own trusted hardware from an inspectable design. Trust
would concern correct execution and protection of that party's secrets,
with fabrication and physical-access assumptions stated. Any reliance on a
peer's hardware would require separate justification.

Preventing unilateral early learning followed by strategic withdrawal is a
research objective. For each selected task, we seek a mechanism that prevents
one party from obtaining the protected result while withholding the messages
needed for the other party's agreed output. The analysis must establish what
protection holds when a participant stops or delays communication and which
communication and hardware assumptions it requires. This output-release
mechanism remains to be constructed and demonstrated.

Correct evaluation does not establish input truthfulness or beneficial
incentives. Any assumption of rational behavior requires an explicit utility
model.

\section{A Small, Falsifiable Starting Point}
The current technical anchor is a finite-game question: can a party
demonstrate that a committed, deterministic Tic-Tac-Toe policy never loses
from the empty board, for a fixed role, without disclosing the policy?
The predicate quantifies over every legal opposing continuation. It does not
assert optimal play at every state, ownership of the strategy, or adherence
to it in later play. Tic-Tac-Toe is a bounded testbed with a checkable
state space; its policies are not presumed to be valuable secrets. It tests
whether a claim about a private decision rule can be carried consistently
through a mathematical relation and an executable representation before
considering more consequential decisions.

The intended proof relation illustrates the composition issue. Let $p$
contain the game version, role, initial state, and session context; let $s$
be a canonically encoded policy and $r$ commitment randomness. For public
commitment $c$, the proposed relation is
\begin{equation}
 R((p,c);(s,r)) =
 [c=\operatorname{Com}(p,s;r)]\ \land\ V(p,s),
\end{equation}
where $V$ checks policy validity and non-loss under the fixed game model.
A prospective zero-knowledge argument of knowledge must bind both checks
to the \emph{same} hidden policy. Separate checks on unrelated witnesses
would not establish this statement. Neither the relation nor a commitment
alone establishes when a policy was fixed; chronology requires an
authenticated ordering mechanism if a later application needs it.

The working repository contains a public policy checker, a fixed directed
acyclic graph evaluator, Boolean circuits for the policy predicate, and
disclosed commitment-opening checks. A Rock--Paper--Scissors demonstration
implements commit--reveal in participant-local state machines, with adversarial
tests that exhibit selective abort. The present artifacts are educational: no private
policy proof, distributed deployment, or hardware implementation is
established. Differential and negative tests check executable behavior;
they do not prove cryptographic hardness or equivalence for every policy.
The next step is to select and jointly justify a compatible commitment
profile and proof representation for the complete relation, then measure
its actual cost and leakage boundary.

\balance
\section{From Specification to Hardware}
The longer-term agenda follows selected components from a formal
specification through explicit syntax or circuit representations and C
reference code to synthesizable register-transfer-level descriptions
(Verilog or VHDL), then FPGA evaluation. High-level synthesis is one possible
route. Behavior-preserving synthesis already has research foundations
\cite{herklotz2021}; this project does not yet supply such a toolchain or
proof of translation correctness.

Each translation must state what it preserves: input encoding, arithmetic,
state transitions, outputs, and error behavior. Proof-oriented circuits and
hardware circuits may need different representations and cost models; a
small proof circuit is not automatically an efficient hardware design.
Evaluation will report proof and verification time, communication, and memory
where relevant, alongside hardware area, latency, throughput, and energy
under a documented workload. Comparisons will use the same functional task
and disclose platform and tool versions.

Functional correspondence alone does not establish confidentiality against
timing, power, electromagnetic observations, or injected faults. These need
separate models and evaluation. Each participant would fabricate its own
hardware to execute the selected protocol components. Evaluation must connect
the inspected design to the fabricated device, stating assumptions about
fabrication, physical access, and device integrity. Synthesis is one stage
of that investigation.

\section{Resources and Support}
We are seeking financial support, research collaborators, and access to
computing and hardware prototyping resources to develop and demonstrate
refereeless cryptographic protocols. Funding would support dedicated
research time, independent cryptographic review, and implementation work.
Access to FPGA boards and measurement equipment would enable the design
and evaluation of hardware for participating parties.

We welcome collaboration in cryptography, game theory, formal methods, and
hardware design. Our initial objective is one carefully specified and
independently reviewed demonstrator, building on the finite-game work
described here. Support would help turn this starting point into inspectable
protocol specifications, reference implementations, adversarial tests, and
measured software and hardware experiments, with the proof status and limits
of each result stated openly.

Prospective supporters and collaborators can reach us at
\href{mailto:contact@waajacu.com}{contact@waajacu.com} to discuss funding,
equipment access, or technical collaboration. The purpose is to make
cooperation possible without requiring parties to surrender their legitimate
secrets or strategic autonomy. Peace is the ethical motivation; it is not a
property a cryptographic proof can establish.

\bibliographystyle{IEEEtran}
\bibliography{references}
\end{document}
